% Task-agnostic benchmark protocol paragraph (main body, 4-6 sentences)
% Insert in Section 6 / Benchmark task specification

\paragraph{Task-agnostic evaluation protocol.}
To ensure the benchmark admits methods beyond \ADS, we define three
formalism-independent tasks with standardized JSON input/output formats.
\textbf{Task~A} (\emph{Multi-Stakeholder Course Ranking}) asks any method to
produce a ranked course set given course texts and stakeholder corpora; we
evaluate coverage, diversity, objective balance, and alignment with a
reference Pareto set (precision, recall, F1).
\textbf{Task~B} (\emph{Stakeholder Trade-off Discovery}) asks for a set of
courses representing distinct trade-offs; we evaluate hypervolume, spacing,
spread, and the fraction of submitted courses on the true Pareto front
(``Pareto precision'').
\textbf{Task~C} (\emph{Robust Recommendation}) requires submissions under
multiple perturbation conditions (encoder noise, weight jitter) and evaluates
bootstrap Jaccard stability, hypervolume retention, and rank correlation across
runs.
A self-contained evaluation script
(\texttt{benchmark\_eval.py}) accepts any method's output---including non-ADS
baselines such as TF-IDF + weighted sum or LLM-based rankers---and computes all
metrics against the released ground-truth scores without requiring access to ADS
internals.

% ---------- Appendix-length version ----------
% Uncomment for supplementary material:
%
% \subsection{Task-Agnostic Benchmark Protocol (Extended)}
% \label{sec:supp-benchmark-protocol}
%
% \paragraph{Motivation.}
% A benchmark coupled to a single formalism limits external
% participation and reproducibility. We decouple the evaluation from
% \ADS by specifying three tasks purely in terms of observable
% inputs (course text artifacts, stakeholder target corpora) and
% outputs (ranked course IDs with optional per-course score vectors).
% Any method---embedding-based, feature-engineered, or
% LLM-based---can produce a conforming submission without importing
% ADS code.
%
% \paragraph{Task A: Multi-Stakeholder Course Ranking.}
% \textit{Input:} Course text corpus $\mathcal{C}$, stakeholder target
% corpora $\{T_k\}_{k=1}^{m}$ (e.g., O*NET occupations for market,
% mission statements for mission).
% \textit{Output:} Ordered list of course IDs
% $[c_1, c_2, \ldots, c_K]$ with optional per-course score vectors.
% \textit{Metrics:}
% (1)~\emph{Coverage}: fraction of objective-space grid cells occupied
% by the selected set;
% (2)~\emph{Diversity}: mean pairwise Euclidean distance in objective
% space;
% (3)~\emph{Objective balance}: normalized Shannon entropy of mean
% objective proportions;
% (4)~\emph{External alignment}: precision, recall, F1, and Jaccard
% against the reference Pareto set.
%
% \paragraph{Task B: Stakeholder Trade-off Discovery.}
% \textit{Input:} Same as Task~A.
% \textit{Output:} Set of course IDs representing distinct
% stakeholder trade-offs.
% \textit{Metrics:}
% (1)~\emph{Hypervolume} (HV): volume of objective space dominated by
% the submitted set, reference point $\mathbf{r} = \mathbf{0}$;
% (2)~\emph{Spacing}: standard deviation of nearest-neighbor distances
% (lower = more uniform);
% (3)~\emph{Spread}: bounding-box volume;
% (4)~\emph{Pareto precision}: fraction of submitted courses on the
% true Pareto front;
% (5)~\emph{Dominance resistance}: fraction of submitted points that
% are non-dominated within the submission.
%
% \paragraph{Task C: Robust Recommendation.}
% \textit{Input:} Same as Task~A, plus a perturbation protocol
% (encoder noise $\sigma$, weight jitter scale).
% \textit{Output:} Primary submission plus $\geq 2$ perturbation-run
% submissions.
% \textit{Metrics:}
% (1)~\emph{Bootstrap Jaccard}: mean pairwise Jaccard similarity of
% course sets across perturbation runs;
% (2)~\emph{HV retention}: $\min(\text{HV}) / \max(\text{HV})$ across
% runs;
% (3)~\emph{Rank stability}: mean pairwise Spearman correlation of
% rankings;
% (4)~\emph{Core stable fraction}: fraction of courses appearing in
% $\geq 80\%$ of runs.
%
% \paragraph{Submission format.}
% Submissions are JSON files with fields:
% \texttt{method\_name} (string), \texttt{task} (``A''/``B''/``C''/``all''),
% \texttt{ranked\_courses} (list of course IDs),
% \texttt{scores} (optional dict of course ID $\to$ objective vector),
% and \texttt{perturbation\_runs} (required for Task~C, list of
% per-run ranked course lists).
% The evaluation script (\texttt{benchmark\_eval.py}) depends only on
% NumPy, SciPy, and the Python standard library.
